\section{Worked examples}
\label{sec:examples}

The examples below show how diverse arguments enter the same normal form.  They are
chosen to exercise exact descent, approximate compactness, relational ambiguity, and a
conditional deep theorem.

\subsection{Euclidean algorithm as a majorized trace}

Input: naturals \(a,b\).  State:
\[
  s_t=(x_t,y_t),\qquad s_0=(a,b).
\]
Transition:
\[
  (x,y)\mapsto
  \begin{cases}
    (x,0),&y=0,\\
    (y,x\bmod y),&y>0.
  \end{cases}
\]
Rank: \(\rank(x,y)=y\).  Clock: \(b+1\).  Invariant:

\begin{enumerate}[label=(\roman*),leftmargin=2.4em]
\item \(\gcd(x_t,y_t)=\gcd(a,b)\);
\item if \(y_t>0\), then \(x_t\bmod y_t<y_t\);
\item once \(y_t=0\), the state is frozen.
\end{enumerate}

The equality of gcds is expressed through the checked divisibility specification or
through preservation of common divisors.  The rank theorem gives termination by the
clock.  The terminal state satisfies \(x_t=\gcd(a,b)\).  Adding coefficient fields and
a parity bit yields the extended Euclidean algorithm and B\'ezout identity.  The
reference test on \((1071,462)\) reaches \(21\) and confirms that the trace length is
bounded by the initial second component plus the frozen step.

This example illustrates that a growing proof state need not have a decreasing total
size.  Only one certified rank field controls termination.

\subsection{Approximate minimum on a rational interval}

Let \(f:[-1,1]\to\mathbb R\) be represented by a sampler with evaluation error
\(\eta(k)\) and a Lipschitz bound \(L\).  Use the dyadic net
\[
   N_k=\{-1+i2^{-k}:0\le i\le2^{k+1}\},
\]
with radius \(\delta(k)=2^{-k}\) and modulus
\(\omega(k)=L2^{-k}\).  Evaluate every point at error \(\eta(k)\), choose the least
sample minimum, and apply \cref{thm:finite-net-min}.  The selected point \(p_k\)
satisfies
\[
   f(p_k)\le f(x)+L2^{-k}+2\eta(k)
   \quad\text{for every }x\in[-1,1].
\]

For \(f(x)=(x-1/3)^2\), the exact reference implementation on the quarter-grid chooses
\(1/4\).  An object proof need not mention the real infimum.  It proves the finite
comparison against any rational input carrying a nearest-grid witness; the semantic
corollary follows from the sampler interpretation.

If one additionally proves strong convexity, the approximate minimizer clouds have a
shrinking diameter, and \cref{thm:shrinking-viability} yields a coherent name for the
unique minimizer.  Without uniqueness, only the value and approximate cloud are
canonical.

\subsection{Polynomial roots through a collision}

Consider a monic polynomial family \(P_t(z)\) and a finite root routine that returns
rational disks covering all complex roots with multiplicity.  At a parameter where two
roots collide, no continuous exact ordering of the individual roots is expected.  The
\(\FRNF\) presentation stores:

\begin{itemize}
\item the disk cloud and multiplicity/count certificates;
\item compatibility edges between overlapping refined disks;
\item a movement bound between refinements; and
\item coefficient reconstruction bounds for the product of the disk centers.
\end{itemize}

Backward viability chooses finite compatible refinements.  The elementary symmetric
functions of any valid root multiset are stable because they reconstruct the
coefficients.  Thus the theorem about the factorization identity is exact even though
internal labels are not.  If a disk is isolated from all others by a rational gap,
\cref{thm:gap-collapse} recovers an exact branch label for that isolated cluster.

The lesson is that collision is not a failure of the root theorem.  It is a failure of
an overstrong output type.

\subsection{Variable-degree deflation with dual horizons}

Suppose a fixed coefficient evaluator gives approximations to a monic polynomial of
numerical degree \(n\).  A one-root routine at output precision \(t\) reads coefficient
rows through \(U(n,M,t)\).  Deflating by the approximate root shifts coefficient
precision and the residual proof requires comparisons through a larger bound
\(V(n,M,t)\).  Repeating the routine \(n\) times creates a changing parameter chain.

Compile it as follows.

\begin{enumerate}[label=\arabic*.,leftmargin=2.2em]
\item Define the stage state: remaining degree, normalized coefficient table, roots
already found, and error budget.
\item Compute all later verification demands backward from the final coefficient
identity.  At each stage propagate \(V\), not only \(U\).
\item Materialize the original coefficient name only through the resulting maximum.
\item Run deflation forward for exactly \(n\) stages.
\item Prove by QF induction a telescoping coefficient identity and accumulated error
bound.
\end{enumerate}

The degree is ordinary data, so this is one primitive recursion uniform in \(n\).  The
coefficient evaluator is statically fixed, respecting the no-universal-evaluator
barrier.

\subsection{Stable simplicial maps}

Let a represented map between compact spaces induce, at finite observation levels, a
carrier cloud: each source vertex maps to a nonempty set of target vertices contained
in one simplex.  Any selection gives a simplicial map.  Any two selections are
contiguous, so the explicit prism operator supplies a chain homotopy.  Therefore the
induced homology map is independent of tie-breaking.

This is a stable observer in algebraic form.  The internal vertex choices are
relational; the homology map is exact and branch-independent.  A proof requires only
finite complexes, bounded matrix operations over a finite ring, the carrier
certificate, and the prism identity.  No exact continuous selector is needed.

\subsection{Height bound implies complete Diophantine list}

Let \(F_p(\vec x)=0\) be a polynomial family and suppose a boundary formula supplies
\[
  F_p(\vec x)=0\quad\Longrightarrow\quad H(\vec x)\le B(p).
\]
Enumerate every canonical rational/integer tuple of height at most \(B(p)\), evaluate
\(F_p\) exactly, and retain the zeros.  Soundness is immediate from exact evaluation.
Completeness follows by the height formula and coverage of the enumeration.  All loops
are bounded by a primitive-recursive function of \(p\) and \(B(p)\).

The resulting theorem is
\[
  \operatorname{HeightBound}(p,B(p))
  \longrightarrow
  \operatorname{CompleteZeroList}(p,L(p)).
\]
For a special curve family, a descent theorem may provide the height bound and close
the schedule.  For general Siegel or Faltings, the finite compiler is still complete as
a conditional result, while the deep producer remains visible.
